\documentclass[a4paper,11pt]{article}

\usepackage[english]{babel}
\usepackage[utf8]{inputenc}
\usepackage{geometry}
\usepackage{graphicx}
\usepackage{hyperref}
\usepackage{wrapfig}

\geometry{a4paper,left=2.0cm,right=2.0cm,top=1.5cm,bottom=1.5cm}

\begin{document}
\pagestyle{empty}

\begin{center}
 \huge{\textbf{Curriculum Vitae}}
 \\
 \huge{Álvaro Fernandes de Abreu Justen}
 \\
 \small{(also known as Turicas)}
 \large{\\}
 \large{\today}
 \\
 \small{The latest version of this file is available at:
        \href{http://turicas.info/curriculum}{turicas.info/curriculum}}
\end{center}

\newcommand{\titulo}[1]{\section*{#1}}
\renewcommand{\labelitemi}{$\diamond$}
\renewcommand{\labelitemii}{$\rightarrow$}

\titulo{Introduction}
\begin{wrapfigure}{l}{0.2\textwidth}
	\begin{center}
		\includegraphics[width=0.2\textwidth]{turicas}
	\end{center}
\end{wrapfigure}

Álvaro Justen is a free/\textit{libre} software user, collaborator and activist
since 2004; goes to many conferences as speaker, attendee and organizer as
well; gives programming and electronics courses and develops software using
mainly the Python programming language. Maker and natively curious, he's always
learning about everything, not only on the technology field: organic food
growing and hypnosis, for example.

Today he is a digital nomad: travels through Brazil (and other countries)
disseminating knowledge through talks and courses (specially about Python,
Arduino, free software and open data) while works (remotely) in many projects
related to Web development, open data and data analysis and also as an expert
at HackHands. Loves coffee, to cook, to travel and to ride a bike.


\titulo{Contact \& Social Media}
	\begin{itemize}
		\renewcommand{\labelitemi}{}
		\item E-mail \& Google Hangouts: \url{alvarojusten@gmail.com}
		\item Mobile, Signal, Telegram \& WhatsApp: +55 21 9 9898-0141
		\item Skype, GitHub, Twitter, Instagram \& SlideShare: turicas
		\item LinkedIn \& Facebook: alvarojusten
		\item Blog: \href{http://turicas.info/}{turicas.info}
	\end{itemize}


\titulo{Professional Experience}
	\begin{itemize}
		\item \href{http://generonumero.media}{\textbf{Genero e Número}}: works
			creating software to download, convert and normalize data related
			to gender inequalities in Brazil --- since May 2016.
		\item \href{http://hackhands.com}{\textbf{HackHands}}: works as an
			expert into the platform, helping programmers to solve their
			problems quickly, in real time --- since Februrary 2015.
		\item \href{http://onyo.com}{\textbf{Onyo}}: worked as senior backend
			developer and systems administrator, using Python, Django, Django
			REST Framework, fabric, nginx, uwsgi and other cool technologies to
			serve thousands of mobile users (100\% remotely) --- from September
			2014 to April 2016.
		\item \href{http://www.cursodearduino.com.br/}{\textbf{Curso de
			Arduino}}: works as a teacher, giving classes about robotics,
			combining concepts of electronics and programming, using the
			Arduino platform and others --- since January 2011.
		\item \href{http://emap.fgv.br/}{\textbf{Fundação Getúlio Vargas -
			Escola de Matemática Aplicada}}: worked as principal software
			developer (Python), systems administrator and was maintainer of the
			\href{http://pypln.org/}{PyPLN} project (distributed natural
			language processing -- Portuguese Language) --- from November 2011
			to September 2014.
		\item \href{http://www.intelie.com.br/}{\textbf{Intelie}}: worked as
			Python developer, systems administrator and was responsible for
			installations, migrations and operations of the company's
			datacenter monitoring software into some huge customers such as
			Globo.com and iG --- from January 2010 to May 2011.
		\item \textbf{Peta5}: was co-founder of this startup which was focused
			on products to the Brazilian digital TV (SBTVD) market; worked as
			Scrum Master, Python developer (mainly using web2py) and systems
			administrator --- from June 2008 to January 2010.
		\item \href{http://www.metaconsultoria.com/}{\textbf{Meta Consultoria}}
			(junior company -- Engeneering School, Fluminense Federal
			University): was selected as trainee and made the whole training
			course (decided on not joinning the company) --- 2008.
		\item \textbf{\textit{Freelancer}}: worked as Web developer using HTML,
			CSS, JavaScript, PHP and MySQL --- from 2003 to 2008.
	\end{itemize}


\titulo{Skills}
	\begin{itemize}
		\item Natural Languages:
		\begin{itemize}
			\item Portuguese: native, fluent.
			\item English: good reading, writing and talking.
			\item Spanish: good reading, intermediate writing and talking.
		\end{itemize}
	\item Programming Languages:
		\begin{itemize}
			\item Python (10+ years of experience)
			\item Bash (10+ years of experience)
			\item Erlang \& Elixir (learning)
			\item Elm (learning)
			\item HTML, CSS \& JavaScript (level: intermediate-advanced)
			\item C/C++, PHP, Ruby \& Perl (level: basic)
		\end{itemize}
		\item Some tools (software):
		\begin{itemize}
			\item Git \& GitHub, Mercurial -- version control, software
				repository
			\item vim -- the editor
			\item Django, flask \& web2py -- Python Web frameworks
			\item Jira -- project management
			\item markdown \& \LaTeX\ -- markup languages to create pretty
				documents
			\item SQLite, MySQL, PostgreSQL, MongoDB, Riak -- databases
			\item nginx, uwsgi, SSH, fabric, docker, AWS -- systems
				administration
			\item Debian GNU/Linux
		\end{itemize}
		\item Methodologies \& Other Skills:
		\begin{itemize}
			\item Scrum and eXtreme Programming
			\item Test-driven development
			\item Development and access to REST APIs, Web crawling and parsing
			\item GNU/Linux systems administration
			\item Screencasts
			\item Remote working
			\item Conference organizing
		\end{itemize}
	\end{itemize}


\titulo{Remarkable Achievements}
	\begin{itemize}
		\item \href{http://pythonquito.tk}{PythonQuito} Organizer
			(\textit{Charlas Pythónicas} and \textit{Django Girls}) --- Quito,
			March 2016.
		\item Creation and maintenance of
			\href{https://github.com/turicas/rows}{rows}: a Python library and
			command-line interface (free/libre software, of course), which
			helps a lot when dealing with tabular data extraction,
			normalization and conversion.
		\item Remodeling and development of new versions the
			\href{http://pypln.org/}{PyPLN} project: distributed natural
			language processing (free software) -- from 2013 to 2014.
		\item \href{http://2012.pythonbrasil.org.br/}{PythonBrasil[8]}
			Organizer, the biggest Python conference in Latin America (approx.
			450 attendees from all over the world, 4 days of duration) --- Rio
			de Janeiro, November 2012.
		\item Brazilian reference on Arduino, for a huge number of talks and
			courses on the topic.
		\item Foundation and coordination of the
			\href{https://groups.google.com/forum/#!forum/arduinrio}{ArduInRio}
			user group --- Rio de Janeiro, June 2010.
		\item Foundation and coordination of the
			\href{http://dojorio.org}{Coding Dojo de Niterói} group ---
			Niterói, October 2009.
		\item SeTel ("Telecommunications Week") Organizer, Fluminense Federal
			University (approx. 180 attendees) --- Niterói, October 2007.
	 \end{itemize}


\titulo{Academic Education}
	\begin{itemize}
		\item Incomplete graduation in
			\href{http://telecom.uff.br/}{Telecommunications Engeneering} at
			\href{http://www.uff.br/}{Fluminense Federal University} ---
			from 2005 to 2010.
		\item High School: Colégio Ruy Barbosa, Três Rios/RJ.
		\item Elementary School: Escola Nossa Senhora de Fátima, Três Rios/RJ.
	 \end{itemize}


 \titulo{Academic Highlights -- Fluminense Federal University}
	\begin{itemize}
		\item Scholarship on Scientific Research on ReMoTE project (Network
			Monitoring for Power Transmission Lines),
			\href{http://www.midiacom.uff.br/ }{MídiaCom Laboratory} --- from
			December 2008 to January 2010.
		\item Scholarship on Scientific Reserach on
			\href{http://complex.if.uff.br/}{Complex Systems and Statistical
			Physics Group} (many projects), Physics Institute, guided by
			Professor \href{http://profs.if.uff.br/tjpp/}{Thadeu Penna} ---
			from April 2007 to November 2008.
		\item Volunteer Math Teacher, Pre-College Preparatory Course,
			Engineering School --- from March 2007 to December 2007.
		\item Scholarship on \href{http://pet.telecom.uff.br/}{Tutorial
			Education Program on Telecommunications (PET-Tele)}, acting on
			tasks related to research, teaching and university extension to
			enhance the Telecommunications Engineering course, guided by
			Professor \href{http://www.telecom.uff.br/~delavega/}{Alexandre de
			la Vega} --- from September 2006 to March 2008.
		\item Volunteer as beta-tester on the pilot project of University Mesh
			Networks, \href{http://www.midiacom.uff.br/}{MídiaCom Laboratory}
			(GT-Mesh/RNP) --- 2006.
		\item GNU/Linux systems administrator and Physics experiments guide at
			\href{http://www.uff.br/casadadescoberta/}{Casa da Descoberta}
			(Science Dissemination Center) --- July 2005 to September 2006.
	\end{itemize}

\end{document}
